%; whizzy paragraph -pdf xpdf -latex ./whizzypdfptex.sh
%; whizzy-paragraph "^\\\\begin{frame}"
% latex beamer presentation.
% platex, latex-beamer でコンパイルすることを想定。 

%     Tokyo Debian Meeting resources
%     Copyright (C) 2009 Junichi Uekawa
%     Copyright (C) 2009 Nobuhiro Iwamatsu

%     This program is free software; you can redistribute it and/or modify
%     it under the terms of the GNU General Public License as published by
%     the Free Software Foundation; either version 2 of the License, or
%     (at your option) any later version.

%     This program is distributed in the hope that it will be useful,
%     but WITHOUT ANY WARRANTY; without even the implied warreanty of
%     MERCHANTABILITY or FITNESS FOR A PARTICULAR PURPOSE.  See the
%     GNU General Public License for more details.

%     You should have received a copy of the GNU General Public License
%     along with this program; if not, write to the Free Software
%     Foundation, Inc., 51 Franklin St, Fifth Floor, Boston, MA  02110-1301 USA

\documentclass[cjk,dvipdfmx,12pt]{beamer}
\usetheme{Tokyo}
\usepackage{monthlypresentation}

%  preview (shell-command (concat "evince " (replace-regexp-in-string
%  "tex$" "pdf"(buffer-file-name)) "&")) 
%  presentation (shell-command (concat "xpdf -fullscreen " (replace-regexp-in-string "tex$" "pdf"(buffer-file-name)) "&"))
%  presentation (shell-command (concat "evince " (replace-regexp-in-string "tex$" "pdf"(buffer-file-name)) "&"))

%http://www.naney.org/diki/dk/hyperref.html
%日本語EUC系環境の時
\AtBeginDvi{\special{pdf:tounicode EUC-UCS2}}
%シフトJIS系環境の時
%\AtBeginDvi{\special{pdf:tounicode 90ms-RKSJ-UCS2}}

\title{Debian / Ubuntu ユーザーミートアップ in 札幌 2018.07 \\　\\Debianで自宅にリモート接続 / OpenVPN編}
\subtitle{}
\author{Norimitsu Sugimoto (杉本 典充) \\dictoss@live.jp}
\date{2018-07-06}
\logo{\includegraphics[width=8cm]{image200607/openlogo-light.eps}}

\begin{document}

\frame{\titlepage{}}

%\section{}

\begin{frame}{アジェンダ}
  \begin{itemize}
  \item 自己紹介
  \item VPNについて
  \item OpenVPNの紹介
  \item OpenVPNの設定（サーバ、クライアント）
  \item おわりに
  \item 参考資料
  \end{itemize}
\end{frame}


\begin{frame}{自己紹介}
  \begin{itemize}
  \item Norimitsu Sugimoto (杉本 典充)
  \item dictoss@live.jp
  \item Twitter: @dictoss
  \item 大学生のときからDebianを使っています
  \item 仕事はソフトウェア開発者をやってます
  \end{itemize}
\end{frame}


\emtext{VPNについて}

\begin{frame}[containsverbatim]{VPNの種類}
  \begin{itemize}
  \item IPsec系
    \begin{itemize}
    \item L2TP/IPsec
    \end{itemize}
  \item SSL-VPN系
    \begin{itemize}
    \item SSTP
    \item SoftEther
    \item OpenVPN
    \end{itemize}
  \item 利用を推奨しない過去の技術
    \begin{itemize}
    \item PPTP（セキュリティ強度が低く、突破済み）
    \end{itemize}
  \end{itemize}
\end{frame}


\emtext{OpenVPNの紹介}

\begin{frame}[containsverbatim]{OpenVPNの紹介}
  \begin{itemize}
  \item SSL-VPNを使ったVPN接続を行うソフトウェア
  \item https://openvpn.net/
  \item ライセンスはGPLv2
  \item TLS 1.0〜TLS 1.2による暗号通信をサポート
  \item サーバ側の機能とクライアント側の機能の両方をもつ
  \item デフォルトのポート番号は1194/UDPで、TCPも利用可能
  \end{itemize}
\end{frame}


\begin{frame}[containsverbatim]{DebianにおけるOpenVPN}
  \begin{itemize}
  \item Debian 8
    \begin{itemize}
    \item openvpn-2.3.4、openvpn-2.4.0（jessie-backports）
    \end{itemize}
  \item Debian 9 
    \begin{itemize}
    \item openvpn-2.4.0、openvpn-2.4.4（stretch-backports）
    \end{itemize}
  \item openvpnの2.3系と2.4系の違い
    \begin{itemize}
    \item 2.3.3からTLS 1.2が利用可能。ただし使える暗号はTLS 1.0由来のものに限られる（AES-CBC、RSA、DHE）
    \item 2.4からAES-GCMや、ECDHE、ECDSAなどTLS-1.2で利用可能な高度な暗号をサポート
    \item サーバとクライアントで2.3系と2.4系が混在する場合、古いバージョンがサポートする設定に合わせる必要がある
    \end{itemize}
  \end{itemize}
\end{frame}


\emtext{OpenVPNの設定の解説}

\begin{frame}[containsverbatim]{OpenVPNの設定（サーバ）}
  \begin{itemize}
  \item "apt-get install openvpn" でインストール
  \item 証明書や鍵の生成（easy-rsa 2）
  \begin{itemize}
  \item \url{https://openvpn.net/index.php/open-source/documentation/miscellaneous/77-rsa-key-management.html}
  \end{itemize}
  \item サーバ側の設定のサンプル
    \begin{itemize}
    \item \url{https://github.com/OpenVPN/openvpn/blob/master/sample/sample-config-files/server.conf}
  \end{itemize}
  \item 設定ファイルは"/etc/openvpn/*.conf"を利用。複数のVPN設定を同時実行するマルチインスタンスが可能。例）443/TCPと1194/UDPの両方をlistenする
  \end{itemize}
\end{frame}

\begin{frame}[containsverbatim]{OpenVPNの設定（サーバ側を一部抜粋）}
\begin{commandline}
tls-server
tls-version-min 1.2
port 443
proto tcp
dev tun0
persist-key
persist-tun
push "dhcp-option DNS 192.168.2.1"
push "route 192.168.2.0 255.255.255.0"
client-to-client
client-config-dir /etc/openvpn/ccd
tls-auth /etc/openvpn/easy-rsa/2.0/keys/ta.key 0
  (snip)
\end{commandline}
\end{frame}


\begin{frame}[containsverbatim]{OpenVPNの設定（クライアント）}
  \begin{itemize}
  \item "apt-get install openvpn" でインストール
  \item クライアント側の設定のサンプル
  \begin{itemize}
  \item \url{https://github.com/OpenVPN/openvpn/blob/master/sample/sample-config-files/client.conf}
  \end{itemize}
  \item 設定ファイルは、"/etc/openvpn/client.conf"とすること
  \item systemctl start openvpn で起動
  \end{itemize}
\end{frame}

\begin{frame}[containsverbatim]{OpenVPNの設定（クライアント側を一部抜粋）}
\begin{commandline}
client
tls-client
tls-version-min 1.2
dev tun
proto tcp
remote yourserver-fqdn 443
nobind
ca /etc/openvpn/ca.crt
cert /etc/openvpn/myclient1.crt
key /etc/openvpn/myclient1.key
tls-auth /etc/openvpn/ta.key 1
log /var/log/openvpn.log
log-append /var/log/openvpn.log
\end{commandline}
\end{frame}

\begin{frame}[containsverbatim]{おわりに}
  \begin{itemize}
  \item Debianパッケージで提供しているOpenVPNを説明しました
  \item サーバ側とクライアント側の設定を説明しました
  \item 速度重視の場合は、UDPを使ってください
  \item つながらない、切れやすい場合は、openvpnの設定ファイルのMTU値を少し小さくしてみてください
  \item 北海道の広大な地でも、VPN接続できればノートPCからどこでも自宅のネットワークへ接続できます
  \end{itemize}
\end{frame}


\begin{frame}{参考情報}
  \begin{itemize}
  \item \url{https://wiki.debian.org/OpenVPN}
  \item \url{https://openvpn.net/index.php/open-source/documentation.html}
  \end{itemize}
\end{frame}

\end{document}


;;; Local Variables: ***
;;; outline-regexp: "\\([ 	]*\\\\\\(documentstyle\\|documentclass\\|emtext\\|section\\|begin{frame}\\)\\*?[ 	]*[[{]\\|[]+\\)" ***
;;; End: ***
